\par
\SourceLeaf{117}
Momentum inertiae acus $1$ per experimenta anteriora iam cognitum erat (vid. art. $11$), quae prodiderant $K=4228732400$, millimetro et milligrammate pro unitatibus acceptis.

Variatio thermometri in utroque atrio, ubi apparatus stabiliti sunt, per totum experimentorum tempus tam parva erat, ut superfluum sit eius rationem habere.

Aggrediamur iam calculum horum experimentorum, ut intensitas magnetismi terrestris $T$ inde eruatur. Inaequalitas oscillationum acus $2$ levem istius intensitatis variationem manifestat: quo itaque de valore determinato sermo esse possit, reducemus tempus observatum oscillationis acus $1$ ad statum medium magnetismi terrestris intra secundam observationum partem. Reductionem aliam requirit hoc tempus propter retardationem chronometri, tertiamque propter torsionem fili. Hoc modo prodit tempus unius oscillationis acus $1$ reductum
\[
\begin{aligned}
&=15{,}22450\times\frac{17{,}29484}{17{,}29995}\mathpunct{.}\frac{86400}{86385{,}76}\mathpunct{.}\BareRoot{\frac{598{,}4}{597{,}4}}\\
&=15''23530=t
\end{aligned}
\]

Hinc deducitur valor producti $TM=\frac{\pi\pi K}{tt}=179770600$. Parvula differentia inter hunc valorem atque eum, quem supra art. $11$ pro die $11$. Sept. invenimus, variationi tum magnetismi terrestris tum status magnetici acus tribuenda est.

E deflexionibus observatis obtinemus
\[
F=\frac{R'^5\operatorname{tang}v'-R^5\operatorname{tang}v}{R'R'-RR}=113056200
\]
si millimetrum pro unitate accipimus, atque hinc
\[
\frac MT=\frac12F\left(1+\frac{\theta}{Tm}\right)=56606437
\]
Comparatio huius numeri cum valore ipsius $TM$ tandem producit
\[
T=1{,}782088
\]
tamquam valorem intensitatis vis magneticae terrestris horizontalis die $18$. Septembris hora $5$.

\SourceLeaf{118}
\Article{24}
Experimenta praecedentia facta sunt in observatorio, loco apparatuum ita electo, ut ferrum a vicinia quantum licuit arceretur. Nihilominus dubitari nequit, quin ferri moles, in parietibus, fenestris et ianuis aedificii copiose sparsae, imo etiam partes ferreae instrumentorum astronomicorum maiorum, in quibus per ipsam vim magneticam terrestrem magnetismus elicitur, effectum neutiquam insensibilem in acus suspensas exerceant. Vires hinc oriundae magnetismi terrestris tum directionem tum intensitatem aliquantulum mutant, experimentaque nostra non valorem purum intensitatis magnetismi terrestris, sed valorem pro loco apparatus $A$ modificatum exhibent. Haec modificatio, quamdiu moles ferreae locum non mutant, ipsaque elementa magnetismi terrestris (puta intensitas et directio) non magnopere mutantur, sensibiliter constans manere debet, quae vero ipsius sit quantitas, hactenus quidem ignotum est, attamen vix crediderim, eam ultra unam duasve partes centesimas valoris totalis ascendere. Ceterum haud difficile foret, quantitatem, proxime saltem, per experimenta determinare, observatis oscillationibus simultaneis duarum acuum, quarum altera in observatorio loco sueto, altera subdiu in distantia satis magna ab aedificio aliisve ferri molibus turbantibus suspendenda esset, et quae dein vices suas commutare deberent. Sed hactenus haec experimenta exsequi non vacavit: tutissimum vero remedium afferet aedificium peculiare, observationibus magneticis destinatum, munificentia regia mox exstruendum, a cuius fabrica ferrum omnino excludetur.

\Article{25}
Praeter experimenta allata permulta alia similia exsequuti sumus, etsi, tempore anteriori cura multo laxiori. Iuvabit tamen, quae e singulis prodierunt, hic in unum conspectum producere, omissis iis, quae ante apparatus subtiliores stabilitos, per alia rudiora subsidia in acubus diversissimarum dimensionum prodierunt, etiamsi \textit{omnia} approximationem saltem ad veritatem praebuerunt. Ecce valores ipsius $T$ per repetita experimenta subinde erutos:
\SourceLeaf{119}
\begin{center}
\begin{tabular}{c|l|l}
Numerus&Tempus, $1832$&$T$\\
I&Maii $21$&$1{,}7820$\\
II&Maii $24$&$1{,}7694$\\
III&Iun. $4$&$1{,}7713$\\
IV&Iun. $24$—$28$&$1{,}7625$\\
V&Iul. $23$, $24$&$1{,}7826$\\
VI&Iul. $25$, $26$&$1{,}7845$\\
VII&Sept. $9$&$1{,}7764$\\
VIII&Sept. $18$&$1{,}7821$\\
IX&Sept. $27$&$1{,}7965$\\
X&Octobr. $15$&$1{,}7860$
\end{tabular}
\end{center}
Experimenta V—IX omnia in eodem loco facta sunt, contra I—IV in locis aliis; experimentum X proprie est mixtum, quum deflexiones loco quidem sueto observatae sint, oscillationes vero alio. Experimentis VII et VIII aequalis fere cura impensa est; contra experimentis IV, V, VI, X paullo minor, experimentisque I—III cura multo laxior. In experimentis I—VIII adhibitae sunt acus diversae quidem, sed eiusdem fere ponderis et longitudinis (pondus erat inter $400$ et $440$ grammata); contra experimento X inserviit acus, cuius pondus $1062$ grammatum, longitudo $485$ millimetrorum. Experimentum IX eum tantummodo in finem susceptum est, ut appareret, quemnam praecisionis gradum per acum minusculam attingere liceat: pondus acus adhibitae erat tantummodo $58$ grammatum, ceterum cura haud minor, quam in experimentis VII et VIII. Nullum est dubium, subtilitatem observationum notabiliter auctum iri, si acus adhuc graviores, e. g. quarum pondus ad $2000$ vel $3000$ grammata surgat, in usum vocentur.

\Article{26}
Dum intensitas magnetismi terrestris $T$ per \textit{numerum} $k$ exprimitur, huic subest unitas certa $V$, puta vis cum illa homogenea, cuius nexus cum aliis unitatibus immediate datis in praecedentibus quidem continetur, attamen modo aliquantulum complicatiori: operae itaque pretium erit, hunc nexum hic denuo producere, ut, quam mutationem patiatur numerus $k$, si loco unitatum fundamentalium ab aliis proficiscamur, elementari claritate ob oculos ponatur.

\SourceLeaf{120}
\expandafter\def\csname HistoricalNoteText13\endcsname{Vix necesse erit monere, significationes literis antea tributas hic cessare.}
Ad stabiliendam unitatem $V$ proficisci oportuit ab unitate magnetismi liberi $M$\HistoricalNote{13} atque unitate distantiae $R$, statuimusque $V$ aequalem vi ipsius $M$ in distantia $R$.

Pro unitate $M$ adoptavimus eam quantitatem fluidi magnetici, quae in quantitatem aequalem $M$ in distantia $R$ collocatam agens producit vim motricem (aut si mavis pressionem) aequalem ei $W$, quae pro unitate accipitur, i. e. aequalem vi, quam exercet vis acceleratrix $A$ pro unitate accepta in massam $P$ pro unitate acceptam.

Ad stabiliendam unitatem $A$ duplex via patet: scilicet vel depromi potest a vi simili immediate data, e. g. a gravitate in loco observationis, vel ab ipsius effectu in corporibus movendis. In modo posteriori, quem in calculis nostris sequuti sumus, duae novae unitates requiruntur, puta unitas temporis $S$ atque unitas celeritatis $C$, ut pro unitate $A$ accipiatur vis acceleratrix ea, quae per tempus $S$ agens producit velocitatem $C$: denique pro hac ipsa accipitur ea, quae motui uniformi per spatium $R$ intra tempus $S$ respondet.

Ita patet, unitatem $V$ a tribus unitatibus vel $R,P,A$ vel $R,P,S$ pendere.

Supponamus iam, loco unitatum $V,R,M,W,A,P,C,S$ alias accipi $V',\allowbreak R',\allowbreak M',\allowbreak W',\allowbreak A',\allowbreak P'.\allowbreak\ C',\allowbreak S'$ simili quo priores modo inter se nexas, atque utendo mensura $V'$ magnetismum terrestrem per numerum $k'$ exprimi, qui quomodo se habeat ad $k$ inquirendum est.

Statuendo
\[
\begin{aligned}
V&=vV'\\
R&=rR'\\
M&=mM'\\
W&=wW'\\
A&=aA'\\
P&=pP'\\
C&=cC'\\
S&=sS'
\end{aligned}
\]
erunt $v,r,m,w,a,p,c,s$ numeri abstracti, atque
\SourceLeaf{121}
\[
\begin{aligned}
kV&=k'V'\quad\text{sive}\quad kv=k'\\
v&=\frac m{rr}\\
\frac{mm}{rr}&=w=pa\\
a&=\frac cs\\
c&=\frac rs
\end{aligned}
\]
e quarum aequationum combinatione obtinemus
\[
\begin{aligned}
\text{I.}\qquad k'&=k\BareRoot{\frac p{rss}}\\
\text{II.}\qquad k'&=k\BareRoot{\frac{pa}{rr}}
\end{aligned}
\]
Quamdiu modum, quem in calculis nostris sequuti sumus, retinemus, formula priori uti oportet; e. g. si loco millimetri et milligrammatis metrum et gramma pro unitatibus accipimus, erit $r=\frac1{1000}$, $p=\frac1{1000}$, adeoque $k'=k$; si lineam Parisiensem et libram Berolinensem, habebimus $r=\frac1{2{,}255829}$, $p=\frac1{467711{,}4}$, adeoque $k'=0{,}002196161k$, unde e. g. experimenta VIII producunt valorem $T=0{,}0039131$.

Si modum alterum sequi, atque gravitatem pro unitate virium acceleratricium adoptare malumus, statuemus pro observatorio Gottingensi $a=\frac1{9811{,}63}$, unde, quamdiu millimetrum et milligramma retinemus, numeri $k$ per $0{,}01009554$ multiplicandi, mutationesque illarum unitatum secundum formulam II tractandae erunt.

\Article{27}
Intensitas vis magneticae terrestris horizontalis $T$, ut ad absolutam reducatur, per secantem inclinationis multiplicanda est. Hanc Gottingae variabilem esse, nostrisque temporibus diminutionem pati, docuerunt observationes ill. \textsc{Humboldt}, qui mense Decembri $1805$ invenit $69^0\ 29'$, mense Septembri $1826$ autem $68^0\ 29'\ 26''$. Equidem d. $23$. Iunii $1832$ adiumento eiusdem inclinatori, quo olim usus erat b. \textsc{Mayer}, inveni $68^0\ 22'\ 52''$, quod retardationem diminutionis indicare videtur, attamen huic observationi minorem fidem haberem, tum propter instrumentum minus perfectum, tum quod observatio in observatorio facta a turbatione molium ferrearum non satis tuta est. Ceterum huic quoque elemento plenior cura in posterum dicabitur.
